\subsection{Structure of Noetherian and complete local rings}

Let A and C be complete Noetherian local rings, and let u be a local homomorphism from C to A, inducing by passage to quotients an isomorphism of \(\kappa_{C}\) onto \(\kappa_{A}\) . Let \((p_{1}, \ldots, p_{m})\) be a sequence generating the ideal \(m_{C}\) of C, and let \(t_{1}, \ldots, t_{n}\) be elements of \(m_{A}\) . Put \(B = C[[T_{1}, \ldots, T_{n}]]\) .

Lemma 3.--- a) There exists a unique homomorphism \(v: \mathbf{B} \to \mathbf{A}\) which extends \(u\) and maps \(\mathbf{T}_i\) to \(t_i\) for \(1 \leqslant i \leqslant n\) .

\begin{enumerate}
\def\labelenumi{\alph{enumi})}
\setcounter{enumi}{1}
\item
  For v to be surjective, it is necessary and sufficient that the sequence \((u(p_{1}), \ldots, u(p_{m}), t_{1}, \ldots, t_{n})\) generates the ideal \(m_{A}\) of A, or equivalently that the classes of these elements modulo \(m_{A}^{2}\) generate \(m_{A}/m_{A}^{2}\) as a vector space over the field \(k_{A}\) .
\item
  For \(v\) to make A a finite B-algebra, it is necessary and sufficient that the sequence \((u(p_1), \ldots, u(p_m), t_1, \ldots, t_n)\) generate an ideal of definition for the \(\mathfrak{m}_{\mathbf{A}}\)-adic topology of A.
\end{enumerate}

Let n denote the ideal of the ring B generated by \(T_{1}, \ldots, T_{n}\) . Every homomorphism v from B to A which extends u and such that \(v(T_{i}) = t_{i}\) maps n into \(m_{A}\) , hence is continuous when B is equipped with the n-adic topology. The existence and uniqueness of v then follow from A, IV, p. 26, prop. 4.

The ring \(B = C[[T_{1}, \ldots, T_{n}]]\) is a complete Noetherian local ring (III, § 2, n° 10, cor. 6 of th. 2 and n° 6, prop. 6), whose maximal ideal \(m_{B}\) is generated by \(p_{1}, \ldots, p_{m}, T_{1}, \ldots, T_{n}\) . We therefore have \(v(m_{B}) \subset m_{A}\) and v defines a homomorphism gr(v) from gr(B) = \(\sum_{n=0}^{\infty} m_{B}^{n}/m_{B}^{n+1}\) to gr(A) = \(\sum_{n=0}^{\infty} m_{A}^{n}/m_{A}^{n+1}\) . Now the ring gr(A) is generated by \(A/m_{A} = \kappa_{A}\) and \(m_{A}/m_{A}^{2}\) ; gr(v) induces an isomorphism of \(\kappa_{B} = \kappa_{C}\) onto \(\kappa_{A}\) , and the classes modulo \(m_{B}^{2}\) of the elements \(p_{1}, \ldots, p_{m}, T_{1}, \ldots, T_{n}\) generate \(m_{B}/m_{B}^{2}\) as a vector space over \(\kappa_{B}\) ; moreover, v is surjective if and only if gr(v) is surjective (III, § 2, n° 8, cor. 2 of th. 1). This proves b).

The ideal of A generated by the sequence \((u(p_{1}), \ldots, u(p_{m}), t_{1}, \ldots, t_{n})\) is none other than \(v(\mathfrak{m}_{\mathrm{B}})\) A. Since \(m_{A}\) contains \(v(\mathfrak{m}_{\mathrm{B}})\) , A is a Zariski ring for the \(v(\mathfrak{m}_{\mathrm{B}})\) A-adic topology. The ring \(A/v(\mathfrak{m}_{\mathrm{B}})\) A is Artinian if and only if its length as an A-module is finite. But since every simple module over A is annihilated by \(m_{A}\) and since, by hypothesis, \(A/m_{A}\) and \(B/m_{B}\) are isomorphic, this occurs if and only if the dimension over the field \(B/m_{B}\) of the vector space \(A/v(\mathfrak{m}_{\mathrm{B}})\) A is finite. By IV, § 2, n° 5, cor. 2 of prop. 9, one therefore sees that \(v(\mathfrak{m}_{\mathrm{B}})\) A is a definition ideal of A if and only if the dimension of \(A/v(\mathfrak{m}_{\mathrm{B}})\) A over \(B/m_{B}\) is finite. This is indeed the case if A is a finite B-algebra.

Suppose that \(v(\mathfrak{m}_{\mathrm{B}})\) A is a definition ideal of A. The \(\mathfrak{m}_{\mathrm{B}}\)-adic topology of the B-module A then coincides with the \(\mathfrak{m}_{\mathrm{A}}\)-adic topology of the ring A, hence is separated. Since \(\mathrm{A} / v(\mathfrak{m}_{\mathrm{B}})\) A is a finitely generated module over \(\mathrm{B} / \mathfrak{m}_{\mathrm{B}}\) , A is a finitely generated B-module (III, § 2, n° 3, example 3 and n° 9, cor. 1 of prop. 12). This proves c).

Lemma 4. --- Suppose that the Noetherian local ring C is regular, and that \((p_1, \ldots, p_m)\) be a coordinate system of C (VIII, § 5, no. 1, def. 1).

\begin{enumerate}
\def\labelenumi{\alph{enumi})}
\item
  If the sequence \((u(p_1), \ldots, u(p_m), t_1, \ldots, t_n)\) is a secant for A (VIII, § 3, no. 2, def. 1), the homomorphism \(v: \mathbf{B} \to \mathbf{A}\) is injective.
\item
  For \(v\) to be injective and make A a finite algebra over B, it is necessary and sufficient that \((u(p_1), \ldots, u(p_m), t_1, \ldots, t_n)\) be a maximal secant sequence for A. Then A has dimension \(m + n\) .
\end{enumerate}

For the sequence \((u(p_1), \ldots, u(p_m), t_1, \ldots, t_n)\) to be a maximal secant sequence for A, it is necessary and sufficient that it generate an ideal of definition of A, and that A have dimension \(m + n\) (VIII, § 3, no. 2, Th. 1). By Lemma 3, c), this is equivalent to saying that A is a finite B-algebra and a ring of dimension \(m + n\) . Now C is an integral Noetherian ring of dimension \(m\) , hence \(\mathbf{B} = \mathbf{C}[[\mathbf{T}_1, \ldots, \mathbf{T}_n]]\) is an integral Noetherian ring of dimension \(m + n\) (VIII, § 3, no. 4, cor. 3 of prop. 8). If A is a finite B-algebra, and if \(\mathfrak{a}\) is the kernel of \(v\) , one has \(\dim(\mathbf{A}) = \dim(\mathbf{B}/\mathfrak{a})\) (VIII, § 2, no. 3, Th. 1, c)); since B is an integral ring of finite dimension, we have \(\dim(\mathbf{B}/\mathfrak{a}) < \dim(\mathbf{B})\) if \(\mathfrak{a} \neq \{0\}\) (VIII, § 1, no. 3, prop. 6, e)). Therefore, if A is a finite B-algebra, \(v\) is injective if and only if A has dimension \(m + n\) . This proves b).

Suppose that the sequence \((u(p_1), \ldots, u(p_m), t_1, \ldots, t_n)\) of elements of \(\mathfrak{m}_{\mathbf{A}}\) is secant. One can adjoin to it (VIII, § 3, no. 2, Th. 1) elements \(t_{n+1}, \ldots, t_{n+r}\) of \(\mathfrak{m}_{\mathbf{A}}\) to make it a maximal secant sequence. By the preceding, there then exists an injective homomorphism \(w\) of \(\mathrm{C}[[\mathrm{T}_1, \ldots, \mathrm{T}_n, \mathrm{T}_{n+1}, \ldots, \mathrm{T}_{n+r}]] = \mathrm{B}[[\mathrm{T}_{n+1}, \ldots, \mathrm{T}_{n+r}]]\) which extends \(v\) and maps \(\mathrm{T}_{n+j}\) to \(t_{n+j}\) for \(1 \leqslant j \leqslant r\) . Therefore \(v\) is injective. This proves \(a\) ).

THEOREM 3. --- Let A be a local, Noetherian, complete ring whose residue field k has characteristic p. Let C be a p-ring of infinite length, whose residue field is isomorphic to k (no. 3, prop. 5).

\begin{enumerate}
\def\labelenumi{\alph{enumi})}
\item
  Let m be the dimension of the vector space \(\mathfrak{m}_{\mathbf{A}}/(\mathfrak{m}_{\mathbf{A}}^{2} + p\mathbf{A})\) over the field k. There exists an ideal a of the ring \(\mathrm{C}[[\mathrm{T}_{1}, ..., \mathrm{T}_{m}]]\) such that A is isomorphic to \(\mathrm{C}[[\mathrm{T}_{1}, ..., \mathrm{T}_{m}]]/\mathfrak{a}\) .
\item
  Let d be the dimension of A. Suppose that \(p1_{A}\) is not a zero divisor in A. Then there exists a subring A' of A isomorphic to \(C[[T_{1}, \ldots, T_{d-1}]]\) and such that A is a finite algebra over A'.
\end{enumerate}

Let C' be a Cohen subring of A (no. 2, th. 1). Since C has infinite length, there exists a homomorphism from C onto C' (no. 3, prop. 4). Consequently, there exists a local homomorphism \(u:C\to A\) . Choose elements \(t_{1},...,t_{m}\) of \(\mathfrak{m}_{\mathbf{A}}\) whose classes form a basis of the vector space \(\mathfrak{m}_{\mathbf{A}}/(\mathfrak{m}_{\mathbf{A}}^{2}+p\mathbf{A})\) over the field \(k\) . We have \(u(p1_{\mathbf{C}})=p1_{\mathbf{A}}\) and Lemma 3, \(b)\) proves the existence of a surjective homomorphism from \(\mathrm{C}[[\mathrm{T}_{1},..., \mathrm{T}_{m}]]\) to A, extending \(u\) and mapping \(\mathrm{T}_{i}\) to \(t_{i}\) for \(1\leqslant i\leqslant m\) . This proves \(a)\) .

Suppose that \(p1_{\mathbf{A}}\) is not a divisor of 0 in A, hence secant for A (VIII, § 3, no. 2, prop. 3). There then exist (VIII, § 3, no. 2, th. 1) elements \(t_1, \ldots, t_{d-1}\) of \(\mathfrak{m}_{\mathbf{A}}\) such that the sequence \((p1_{\mathbf{A}}, t_1, \ldots, t_{d-1})\) is maximal secant for A. The Noetherian local ring C is regular, and \((p1_{\mathbf{C}})\) is a system of coordinates of C. Assertion \(b)\) of th. 3 then follows from lemma 4, \(b)\) .

